\section{Introduction}

GUI visual search is not only a localization problem. In real interfaces, users often search for controls, settings, links, or actions that may not be present on the current screen. A model that is forced to always return a target can therefore hallucinate plausible UI elements and create unsafe or misleading interaction traces.

This draft studies that gap through the SeekUI pipeline. The central question is whether scanpath and UI evidence can help a target-present model make a stopping decision: can the system decide that the target is absent instead of grounding every query to some UI element?

We therefore use the following research question:

\begin{quote}
Can scanpath and UI evidence help GUI visual search models decide when to stop and report target absence, instead of forcing every query to ground to a UI element?
\end{quote}

The core claim is simple: target-absent GUI search benefits from explicit evidence accumulation. We test three forms of evidence: scanpath-based cognitive stopping, OCR candidate matching, and a VLM verifier that sees the screenshot together with path/OCR evidence.

The current draft assets support three empirical points. First, prompt-only SeekUI shows a forced-choice failure mode on target-absent trials. Second, a combined path+OCR verifier substantially reduces absent false-present errors on a held-out image split. Third, full-benchmark VLM comparisons suggest that screenshot reasoning and scanpath/OCR evidence are complementary, with the same combined-verifier trend appearing in a small realistic absent validation set.
